\chapter[Reactive transport with PHT3D]{Reactive transport with PHT3D}
\def\shorttitle{Reactive transport with PHT3D} % Abbreviated chapter title for right top header

In this chapter we introduce coupled modelling of mass transport and geochemical reactions.
While different modelling approaches for the simulation of such reactive transport 
problems exist, we will focus our discussion on the PHT3D model, which
combines the two existing models MT3DMS and PHREEQC-2 through 
a split-operator technique \citep[]{herzer89,phdzysset}.
This technique has more or less become the standard method for solving such combined
physico-biogeochemical problems. The technique involves separating
the processes (e.g., flow, transport of individual chemical
components/species, chemical reactions, microbial activity) within
the numerical model and solving each sub-model independently. The
technique enables use of available, well tested, transport, flow
and reaction packages, as well as associated GUIs for setting up
problems. Within the PHT3D simulator MT3DMS and PHREEQC-2 solve two 
different sets of equations: one for hydrological solute transport and 
one for the chemical reactions. MT3DMS is used to solve the solute 
transport part and PHREEQC-2 computes the concentration changes
that are caused by geocchemical reactions. 

\section{Solute transport}

The reactive transport equation (solved by MT3DMS) for the
$i^{th}$ (mobile) aqueous component is (in indical notation):
\begin{equation}
{\frac{\partial C_i} {\partial t} = \frac{\partial} {\partial
x_{\alpha}}
 (D_{\alpha\beta} \frac{\partial C_i} {\partial x_{\beta}}) -
 \frac{\partial} {\partial x_{\alpha}} (v_{\alpha} C_i) + \frac{q_s C^s_i}{\theta} + r_{reac,i}}
\label{eq:ad_reac_mobile}
\end{equation}
and for immobile entities, e.g., minerals:
\begin{equation}
{\frac{\partial C_i} {\partial t} = r_{reac,i}}
\label{eq:ad_reac_immobile}
\end{equation}

where $v_{\alpha}$ is the pore-water velocity in direction
$x_{\alpha}$, $D_{\alpha\beta}$ is the hydrodynamic dispersion
coefficient tensor, $q_s$ is a volumetric flow rate per unit
volume of aquifer representing fluid sources (positive) and sinks
(negative), $\theta$ is the porosity of the subsurface medium,
$C^s_i$ is the concentration of the source or sink flux,
$r_{reac,i}$ is a source/sink rate due to chemical reaction and
$C_i$ is the total aqueous component concentration of the $i^{th}$
component \cite[]{yehtri89,enges92}. There $C_i$ is defined as:
\begin{equation}s
{C_i = c_i + \sum_{j=1,n_s} \ Y_{j}^s s_j}  
\label{eq:cj}
\end{equation}
where $c_i$ is the molar concentration of the (uncomplexed)
$i^{th}$ aqueous component, $n_s$ is the number of species in
dissolved form that have complexed with the aqueous component,
$Y_j^s$  is the stoichiometric coefficient of the aqueous
component in the $j^{th}$ complexed species and $s_j$ is the molar
concentration of the $j^{th}$ complexed species.

In the current PHT3D version the (local) redox-state, $p$e, is
modeled by transporting chemicals/components in different redox
states separately, while the $p$H is modeled from the (local)
charge balance. Using a sequential operator-splitting technique
\cite[]{herzer89,val92,miller93,phdzysset,zysset94_2,morka95_1,morka95_2,
barry96,barry97,steefel96} the advection and dispersion terms
within (\ref{eq:ad_reac_mobile}) are solved by the transport
module MT3DMS for $n_{tot}$ components and for each time step of a
temporally discretised problem, with
\begin{equation}
n_{tot} = n_{e,nre} + \sum_{l = 1,n_{e,re}}{n_{rs,l}}  \
\label{eq:ntot_manual}
\end{equation}
where $n_{e,nre}$ is the number of (mobile) chemical elements
occurring in only one redox state, $n_{e,re}$ is the number of
elements occurring in multiple redox states and $n_{rs,l}$ is the
appropriate number of different possible redox states of the
$i^{th}$ element. In other words, if chemical elements occur in
multiple, i.e., $n_{rs,l}$ different redox states within a
simulation problem, an appropriate, higher number of separate
transport equations needs to be solved in order to accurately
predict the redox state ($p$e). For example, hydrological
transport of iron-complexes would typically require modelling the
total aqueous component concentrations of $Fe(II)$ and $Fe(III)$
complexes separately.

\section{Geochemical Reactions}

In PHT3D all concentration changes of aqueous components and
immobile entities that result from reactive processes are computed
by PHREEQC-2. It is, in contrast to its precursor models PHREEQE
\cite[]{parkhurst80} and PHREEQC \cite[]{parkhurst95}, capable of
simultaneously solving arbitrary, kinetically controlled reactions
in addition to geochemical equilibrium problems. From the full set
of geochemical reactions that can be handled by PHREEQC-2, a
subset, including aqueous complexation, mineral
precipitation/dissolution, and ion exchange has been implemented
into the present version of PHT3D. This subset will be sufficient
for most subsurface reaction/transport simulation problems.
However, some limitations remain, as will be discussed in
\ref{limit}.

\section{Coupling Procedure}
\label{sec:coupling} Several options to couple transport and
reactive processes through split-ope\-ra\-tor methods exist, e.g.,
sequential, alternating and iterative split-operator methods.
Details of the different techniques have been discussed, e.g., by
\cite{herzer89}, \cite{val92}, \cite{miller93}, \cite{phdzysset},
\cite{zysset94_2}, \cite{morka95_1,morka95_2}, \cite{steefel96}
and \cite{barry96,barry97,barry2002}. In PHT3D a slightly modified
version of the 'standard' sequential split-operator method
described by \cite{walt94a} has been implemented. A
time-discretised form of (\ref{eq:ad_reac_mobile}) for reactive
transport of a component, as defined by (\ref{eq:cj}) is:
\begin{equation}
C_i^{k+1} - C_i^k = L(C_i)^{k+1/2} \ \Delta t + R_{reac,i}\
\label{eq:ss_discretinorg}
\end{equation}
where $L(C_i)^{k+1/2}$ ($ML^{-3}$) is the spatial differential
operator (central in time, i.e., time level $k+1/2$) and $C_i^k$
and $C_i^{k+1}$ ($ML^{-3}$) are the total aqueous component
concentrations of the $i^{th}$ component at the old time level $k$
and the new time level $k+1$, respectively, $\Delta t$ is the time
step length, and $R_{reac,i}$ ($ML^{-3}$) is the difference in
concentrations from before and after a reaction step. In the model
presented by \cite{walt94a}, $R_{reac,i}$ represents the relative
mass of a component generated or consumed instantaneously by an
equilibration of the aqueous solution with a mineral assemblage at
the time $t^{equil}$ ($t^k \le t^{equil} \le t^{k+1}$),
compensating for the chemical pertubation arising from advective
and dispersive transport during $\Delta t$. However, in PHT3D
$R_{reac,i}$ does not only represent concentration changes of
aqueous components due to (equilibrium) precipitation/dissolution
reactions, but also due to any other reactive process handled by
PHREEQC-2, such as ion exchange, biodegradation, etc., depending
on which reactions the $i^{th}$ component participates in. The
concentration change $R_{reac,i}$ might occur as a result of
equilibrium, kinetic or mixed equilibrium/kinetic reactions. In
PHT3D, the (MT3DMS) advection and dispersion modules solve the
concentrations $C_i^{trans}$, i.e., the concentrations at the new
time level $k+1$ that result from physical transport only:
\begin{equation}
C_i^{trans} = L(C_i)^{k+1/2} \ \Delta t + C_i^k\ \label{eq:seq_1}
\end{equation}
for the $n_{tot}$ compounds/components. In contrast to
\cite[]{walt94a} and depending on the actual advection scheme
selected for a simulation (FD, TVD, MOC, MMOC or HMOC), the
transport simulator MT3DMS may further sub-divide the user-defined
time step length $\Delta t$ into several transport steps that
fulfil the relevant stability and/or accuracy criteria for
physical transport (Courant number). Thus, sometimes several
subsequent transport steps (without reactions) might be carried
out by MT3DMS to compute $C_i^{trans}$. Once $C_i^{trans}$ is
computed, the second step, which models the reaction(s), will be
carried out in a sequential manner for each grid cell
individually. The computation of the reaction step within a grid
cell is independent from concentrations at neighboring cells. This
then leads to the concentrations $C_i^{k+1}$ at the new time level
$k+1$, i.e.,
\begin{equation}
C_i^{k+1} = C_i^{trans} + R_{reac,i}\ \label{eq:seq_2}
\end{equation}
For the computation of the source/sink term $R_{reac,i}$, the
concentrations $C_i^{trans}$ from after the transport step(s) are
used as initial concentrations for the chemistry calculation step.
The time step length for the reaction simulation will correspond
to the user-defined time step length $\Delta t$ and not to the
automatically selected transport step size, which depends on the
actual advection scheme used for a simulation. For the computation
of kinetic reactions PHREEQC-2 uses a $5^{th}$ order integration
scheme \cite[]{fehlberg69} which, if a specified tolerance is not
met for the error estimate of the integration, will automatically
subdivide the user-defined time step length $\Delta t$. In that
case, the user-defined time step will be numerically integrated
over two or more subintervals, each of which is integrated with
the $5^{th}$ order integration scheme. Through this procedure the
kinetic reactions are assured to be integrated accurately over the
user-defined time interval. Of course, an increased temporal
splitting error will occur compared to a scheme where transport
and reactions steps alternate more often, e.g., at the frequency
dictated by the transport step size.

However, the model execution times, i.e., CPU-times of the present
scheme can be significantly lower, which often is helpful,
particularly in the early, more conceptual stages of a modelling
project, i.e., where (very high) solution accuracy is not a
primary concern. Accuracy can be improved by reducing the size of
$\Delta t$. For a small enough step size the numerical solution
will become independent of $\Delta t$ as a result of a negligible
splitting error.

In order to reduce computational load, a further optional
modification of the 'standard' sequential split-ope\-ra\-tor
tech\-nique was implemented in PHT3D. The modification allows the
reaction step to be temporally omitted in cells where no reactive
changes are expected. This modification exploits the fact that in
many model applications there may exist zones, sometimes large,
within the model domain where chemical gradients are negligible
between neighbouring grid-cells and where the aqueous solution is
in equilibrium with the mineral assemblage and the cation-exchange
sites. In these zones the reaction term $r_{reac}$ computed by
PHREEQC-2 might be zero for significant portions of the total
simulation time and thus the reaction step, in principle, can be
omitted during these periods. A typical case where this applies is
a steadily spreading point source contamination within an
initially, chemically homogeneous, equilibrated multi-dimensional
model domain, where in the early stages of the simulation
geochemical changes occur in a very small portion of the aquifer.
This concerns in particular cases where advection is the dominant
transport mechanism. To decide whether the reaction step in a
particular grid-cell can be omitted, the joint fulfilment of the
two criteria (three-dimensional case):
\begin{equation}
  {\sum_{k = lay-2}^{lay+2}\sum_{l=col-2}^{col+2}\sum_{m=row-2}^{row+2}\sum_{i=1}^{n_{tot}}
\ \ \mid r_{reac_{k,l,m,i}} \mid \ \ < \ \epsilon_{aqu}
\ \ \ \ \ \ \ \ \ \ \ \ \ \ \ \ \ \ \ \ }
\label{eq:phcact_dis}
\end{equation}
and
\begin{equation}
  {{\sum_{k = lay-2}^{lay+2}\sum_{l=col-2}^{col+2}\sum_{m=row-2}^{row+2}
r_{reac_{k,l,m,pH}}} \ < \ \epsilon_{pH} \ \ \ \ \ \ \ \ \ \ \ \ \
\ \ \ \ \ \ \ }
\label{eq:phcact_ph}
\end{equation}
with
\begin{equation*}
  { 3 \leq k \leq lay_{max}-2, \\
    3 \leq l \leq col_{max}-2, \\
    3 \leq m \leq row_{max}-2   }
\end{equation*}
has been found to be an efficient test, where $r_{reac_{k,l,m,i}}$
is the computed reaction term for the $i^{th}$ component in layer
$k$, column $l$ and row $m$, both $\epsilon_{aqu}$ and
$\epsilon_{pH}$ are user-defined accuracy criteria, $lay$, $col$,
and $row$ are the indices of the layer, column and row of the
(tested) grid-cell, respectively, and $lay_{max}$, $col_{max}$ and
$row_{max}$ are the total numbers of layers, columns and rows,
respectively in the model. The criteria (\ref{eq:phcact_dis}) and
(\ref{eq:phcact_ph}) are tested for each grid-cell after a
PHREEQC-2 reaction step. If both are fulfilled for a grid-cell,
the execution of PHREEQC-2 (but not physical transport) is
deactivated in this cell until it is reactivated.

Reactivation occurs automatically when either criterion
(\ref{eq:phcact_dis}) or (\ref{eq:phcact_ph}) is not fulfilled,
i.e., if reactive changes occur in neighbouring grid cells. Note,
that (\ref{eq:phcact_dis}) does not only include the directly
adjacent cells but also cells that are two layers, two columns or
two rows away. Grid-cells that receive external fluxes (e.g.,
through wells, rivers, recharge, \ldots{}) can not be deactivated and,
furthermore, reaction simulations are always carried out for all
grid-cells at the beginning of each new stress period, including
the very first time step of the first stress period. For certain
simulation problems, such as those where the initial aqueous
solution and/or the inflow solution across model boundaries are
not in equilibrium with the mineral assemblage and mineral
reactions are kinetically controlled, $r_{reac}$ might always be
non-zero in most or all grid-cells.

\section{Program design}
The coupling between MT3DMS and PHREEQC-2 was achieved by several
new subroutines that enable the communication between the two
programs. The MT3DMS code, written in Fortran 90, was used as the
driver program, i.e., the main program, whereas the original
PHREEQC-2 main program, written in the C programming language, has
become a subroutine of MT3DMS. The new additional subroutines,
called by MT3DMS, were also written in the C programming language.
Despite the addition of those new subroutines, the original main
model structures of both models, MT3DMS and PHREEQC-2 was largely
maintained. This bears the notable advantage that future, modified
and improved versions of the two underlying models can be quickly
implemented with a modest effort. The disadvantage of this 'loose'
coupling strategy is the loss of some computational efficiency as
a result of occasional increased overhead, e.g., through the
repeated reading and interpretation of the same database file and
through partial data-communication via ASCII-files. Figure
\ref{fig_pht3d_flowchart_a} shows a flowchart of the PHT3D main
program that indicates the main modifications that were made
compared to the original MT3DMS code. As can be seen in the chart,
the PHT3D-specific modifications concern the
\begin{itemize}
\item data preparation where the new subroutine \textbf{MCRP} is called.
\item newly added subroutines \textbf{PREPH}, \textbf{PHMAIN} and \textbf{POSTPH2}, which
preprocess, execute and postprocess, respectively, the PHREEQC-2 reaction step
\item the newly added subroutine \textbf{PHCACT} which decides after each reaction step
for which grid-cells a reaction simulation will be necessary during the next reaction step.
\end{itemize}

\begin{figure}
\begin{center}
\epsfig{file=pht3d_flowchart_d.eps,width =10.cm ,angle=0,clip=}
\end{center}
\caption[Flowchart of the main program of PHT3D]{Flowchart of the main
program of PHT3D, modified from \citep[]{zheng99}.
The PHT3D-specific changes of the original MT3DMS code are marked in yellow.}
\label{fig_pht3d_flowchart_a}
\end{figure}

As part of the initial data preparation procedure, the MT3DMS main
program calls the newly added \textbf{MCRP} subroutine in PHT3D.
During the execution of this subroutine the PHREEQC-2 interface
package file {\color{red} \textbf{pht3d\_ph.dat}} is read and the
input data are prepared for their use in other PHT3D-specific
subroutines. For example, the subroutine reads, sorts and stores
the number and names of entities included in a simulation and
allocates a unique species number to each of those entities. Given
by the structure of MT3DMS, mobile entities will form the first
'block' of species, followed by immobile entities such as
minerals, immobile bacteria or (immobile) NAPL compounds which
form the second 'block'.

In the course of a model simulation, i.e., during each reaction
step, the three subroutines \textbf{PREPH}, \textbf{PHMAIN} and
\textbf{POSTPH2} preprocess, execute and postprocess the PHREEQC-2
simulations that compute the reaction terms. In detail, the
subroutine \textbf{PREPH}
\begin{itemize}
\item reads the concentrations of all entities from the
appropriate MT3DMS arrays and \item generates an ASCII file
{\color{red} \textbf{phinp.dat}} that will subsequently serve as
input file for PHREEQC-2.
\end{itemize}
Following the creation of {\color{red} \textbf{phinp.dat}}, the
execution of the \textbf{PHMAIN} subroutine (a modified version of
the original PHREEQC-2 \textbf{MAIN} program)
\begin{itemize}
\item performs the computation of the reaction step using
{\color{red} \textbf{phinp.dat}} and the database file
{\color{red} \textbf{pht3d\_datab.dat}} as input files and \item
writes the newly computed concentrations to the file
{\color{red} \textbf{phout\_sel.dat}}, an ASCII file produced by
and during the execution of PHREEQC-2. It contains all relevant
simulation results, i.e., the newly computed concentrations of
mobile and immobile chemicals from after the reaction calculation.
\end{itemize}
Subsequently, the subroutine \textbf{POSTPH2} will
\begin{itemize}
\item read the output file {\color{red} \textbf{phout\_sel.dat}}
that contains those newly computed concentrations \item compute
the concentration changes from before and after the reaction
simulation, i.e., the reaction term $r_{reac}$ \item save the
newly computed concentrations in the appropriate MT3DMS data
arrays.
\end{itemize}
The three above steps/subroutines might be carried out only once
for all $n_{grid,max}$ (= $lay_{max} \times col_{max} \times
row_{max}$) grid-cells of the model domain or, alternatively,
repeatedly for several packets of a user-defined maximum number of
grid-cells.

At the end of each reaction step, i.e., once the reaction terms
$r_{reac}$ are known for all $n_{grid,max}$ grid-cells, the
subroutine \textbf{PHCACT} is invoked to check whether the
criteria defined by equations (\ref{eq:phcact_dis}) and
(\ref{eq:phcact_ph}) are fulfilled and appropriate internal flags
are set for each grid-cell. The settings will be used during the
reaction simulation of the next time step.

\section{Model Structure} 
For the development of PHT3D, the original MT3DMS model has been 
used as a basis to incorporate an interface that prepares and executes 
the PHREEQC-2 reaction step. This model interface was kept relatively 
simple which has the notable advantage that future, modified and 
improved versions of not only the transport model, but also 
the geochemical model, can be quickly implemented with a modest effort. 
The model input/output for the existing MT3DMS features has not been changed. 
Compared to a standard MT3DMS simulation, only two additional input files 
are needed, one which contains information such as the names and types 
of chemicals, parameter values, and the stoichiometry of 
kinetic reactions and a second file that contains the definition 
of both equilibrium and kinetic reactions. The syntax of the 
latter file is analogous to that one used by the 
original PHREEQC-2 database file. Figure \ref{fig:data_flow} shows the data flow within a typical PHT3D application.


\begin{figure}[ht]
\centering
%\includegraphics[width=10cm,bb=42 0 450 250]{pht3d_data_flow.png}
\includegraphics[height=10.5cm, angle=270]{data_flow.eps}
\caption[Data flow in a MODFLOW-based PHT3D simulation.]{Data flow in a MODFLOW-based PHT3D simulation.}
\label{fig:data_flow}
\end{figure}


\section{Limitations} \label{limit}

Only a subset, but not all of the geochemical reactions that can
be computed with the original PHREEQC-2 model is currently
supported by PHT3D. Among those unsupported reactive processes is
the surface complexation model (SCM) which, in some cases, might
be essential to simulate the mobility of ions, e.g., of trace
metals, accurately. While the SCM is not included in the present
version, it is planned to be incorporated into future
versions of PHT3D.

As in the original MT3DMS model and other MT3DMS-based package,
PHT3D simulations are carried out on the base of a flow-field
computed beforehand by a separate flow simulator, typically
MODFLOW. Thus, it can't reproduce the potential impact of reactive
processes on the groundwater flow field and the model is not
suitable to predict, for example, the impact of bioclogging or
mineral precipitation on the hydraulic properties of an aquifer.

Note also that the use of the implicit GCG solver is currently not supported.
Furthermore, the simultaneous use of the MT3DMS chemical reaction
package and PHREEQC-2 as reaction simulator was not tested.

\section{Additional literature}

More detailed information about the PHT3D model and several application 
examples can be found in \cite{pht3d} and in \cite{prommer_gw}. Excellent introductions
into the general area of reactive transport are provided by \cite{steefel96} and by \cite{ap2005}. 

